\section{Discussion}

The results confirm that the \Core{} provides a lightweight, modular
self‑regulation layer.  The event bus is sufficiently fast for the
scale of a single knowledge agent; if the agent were to grow to
thousands of concurrent layers, a distributed backend such as Kafka
would be mandatory.  The adaptive thresholds and chaos detector react
quickly enough to prevent runaway dynamics in low‑dimensional
simulations; their effectiveness in high‑dimensional, real‑world
scenarios remains an open question.

A central limitation is the energy model.  While the CPU and GPU power
estimates are reasonably accurate (within 10\,\% of a wall‑plug
measurement), the FPGA and quantum power models are currently
placeholder values derived from datasheets.  The efficiency criterion also introduces an ethical tension: under
stringent carbon constraints, the system may forgo potentially valuable
but energy‑intensive computations, effectively trading cognitive
capacity for ecological responsibility.  This trade‑off is logged and
can be reviewed by human operators. Also, the per‑structure
efficiency metric relies on a heuristic for information value; a more
principled measure, such as the reduction in variational free energy
that the structure achieves, would align better with the \Apeiron{}
framework’s theoretical grounding.

\paragraph{The Energetic Boundary of Autonomy.}
The thermodynamic cost model introduces a fundamental tension: under severe carbon 
constraints, the agent must trade cognitive capacity for ecological responsibility.  
This effectively creates an “energy‑bound consciousness” — the system cannot think 
all thoughts, only those that fit within its energetic means.  
We argue that this limitation is not a flaw but a necessary feature for any 
long‑lived autonomous agent; it forces the system to internalize the true cost of 
computation, much as biological organisms allocate metabolic resources.  
Future work will explore the philosophical implications of this forced frugality.

\paragraph{Stability of the Meta‑Controller.}
A legitimate concern is whether the coupling between AdaptiveThresholds
and ChaosDetector could create a positive‑feedback loop in which rising
complexity relaxes thresholds so far that genuine divergence is missed.
The system contains three safeguards against this scenario.  First, all
thresholds are strictly clamped to $[0.05,0.95]$ (Section~4.2),
ensuring that even at maximal relaxation a minimal degree of sensitivity
remains.  Second, the trend factor $\gamma_t$ in Equation~4.2 acts as a
restoring force: if the filtered complexity estimate begins to fall
while the raw complexity is still rising—a classic late‑stage
divergence signature—$\gamma_t$ tightens thresholds again.  Third, a
\texttt{DANGER} intervention automatically resets the system to a known‑
good checkpoint and forces a full re‑calibration of all thresholds,
breaking any sustained drift.  We have not observed a runaway positive‑
feedback loop in any of the synthetic or semi‑synthetic workloads tested
so far.

Our trust model assumes that the sensors (psutil, NVML) are
uncompromised.  In a safety‑critical deployment, redundant power
sensors and cryptographic attestation of the energy log would be
required.

Future work will integrate the \Core{} with the higher layers of the
\Apeiron{} Framework, particularly Layer 10 (coherence evaluation) and
Layer 15 (morphogenetic ethics).  We also plan to extend the event bus
with a dedicated stream‑processing language for expressing complex
event patterns.